% GENERATED FILE. Edit canonical lessons, then run npm run generate:books.
% canonical-source-hash: fnv1a64:087885552d8aa7a4
% canonical-lessons: IT-C08-ora, IT-C08-mezzogiorno-mezzanotte

\chapter{Telling Time}
\label{ch:it-telling-time}
\hlchaptermodality{8}

\begin{chapteropening}
\textbf{By the end of this chapter:} \emph{I can tell the time in Italian and use mezzogiorno and mezzanotte for the two hours that need no number.}

Say è mezzogiorno and è mezzanotte, break each into mezzo or mezza plus giorno or notte, and explain why the half changes gender between them.
\end{chapteropening}

\section[ora]{ora — the hour, told with \textquotedblleft{}le\textquotedblright{}}

\label{lesson:IT-C08-ora}



\begin{quote}
Italian's word for \textquotedblleft{}hour\textquotedblright{} is \textbf{ora}, straight from Latin \textbf{hōra} (which Rome took from Greek \emph{h\'{\={o}}rā}, \textquotedblleft{}a time of day\textquotedblright{}). Telling the time in Italian has one neat twist: you use the \textbf{feminine article} and let \emph{ore} (\textquotedblleft{}hours\textquotedblright{}) stay implied.
\end{quote}

\begin{sounds}
\begin{itemize}
\raggedright
  \item \texttt{open-o} — \textbf{ora} = \emph{OH-ra}, open \emph{o}, single tapped \emph{r}.
\end{itemize}
\end{sounds}

\begin{cousinweb}
\textbf{ora} $\leftarrow$ Latin \textbf{hōra} $\leftarrow$ Greek \textbf{h\'{\={o}}rā}. It's the closest of the sisters to the Latin: French wore it to \emph{heure}, but Italian barely touched it — \emph{hōra $\to$ ora} (Italian dropped the silent Latin \emph{h-} in spelling, too). Same word as English \emph{hour}, and as German's borrowed \emph{Uhr}.
\end{cousinweb}

\begin{grammarlens}[title={Grammar Lens: \textquotedblleft{}è l'una\textquotedblright{} but \textquotedblleft{}sono le due\textquotedblright{}}]
Here's the twist. Italian tells the time with the \textbf{feminine article} (because \emph{ora/ore} is feminine), and the verb agrees with the number:

\begin{quote}
\textbf{È l'una.} — \textquotedblleft{}It is one o'clock.\textquotedblright{} (singular: \emph{è} \textquotedblleft{}it is\textquotedblright{} + \emph{l'una} = \textquotedblleft{}the one {[}hour{]}\textquotedblright{}) \textbf{Sono le due.} — \textquotedblleft{}It is two o'clock.\textquotedblright{} (plural: \emph{sono} \textquotedblleft{}they are\textquotedblright{} + \emph{le due} = \textquotedblleft{}the two {[}hours{]}\textquotedblright{}) \textbf{Sono le dieci.} — \textquotedblleft{}It is ten o'clock.\textquotedblright{}
\end{quote}

Literally \textquotedblleft{}\textbf{they are} the two {[}hours{]}\textquotedblright{} — Italian counts the hours as feminine things and drops the word \emph{ore} entirely; the article \textbf{le} carries it. Only \textbf{one o'clock} is singular (\emph{è l'una}); from two on, it's \emph{sono le …}.
\end{grammarlens}

\subsection*{Your turn}
Say these aloud:

\begin{itemize}
\raggedright
  \item \textquotedblleft{}ora\textquotedblright{} — \emph{OH-ra}
  \item \textquotedblleft{}è l'una … sono le due, sono le tre, sono le quattro\textquotedblright{}
  \item \textquotedblleft{}ora / heure / Uhr / hour\textquotedblright{} — all Latin \emph{hōra}
\end{itemize}

\begin{tcolorbox}[breakable,colback=teal!4,colframe=teal!35!black,title={Before you move on}]
What word is \emph{ora} from, and its sisters? (Latin \emph{hōra} $\leftarrow$ Greek \emph{h\'{\={o}}rā}; French \emph{heure}, German \emph{Uhr}, English \emph{hour}.) Why \textquotedblleft{}\emph{sono le} due\textquotedblright{} but \textquotedblleft{}\emph{è l'}una\textquotedblright{}? (Plural hours take \emph{sono le …}; one o'clock is singular \emph{è l'una}.) What word does the \emph{le} stand in for? (\emph{ore}, \textquotedblleft{}hours\textquotedblright{} — feminine.) Next: \textbf{mezzogiorno} and \textbf{mezzanotte} — noon and midnight.
\end{tcolorbox}
\section[mezzogiorno, mezzanotte]{mezzogiorno and mezzanotte — Italy's noon and midnight}

\label{lesson:IT-C08-mezzogiorno-mezzanotte}



\begin{quote}
The two unnumbered hours. \textbf{mezzogiorno} (noon) and \textbf{mezzanotte} (midnight) are, like their French cousins, the \textbf{middle} of the day and the night — built from \emph{mezzo} \textquotedblleft{}half/middle\textquotedblright{} plus \emph{giorno} \textquotedblleft{}day\textquotedblright{} and \emph{notte} \textquotedblleft{}night.\textquotedblright{}
\end{quote}

\begin{sounds}
\begin{itemize}
\raggedright
  \item \texttt{double-zz} — \textbf{mezzo} = \emph{MED-dzo}, the doubled \emph{zz} a hard \textquotedblleft{}ts/dz\textquotedblright{} — lean on it.
\end{itemize}
\end{sounds}

\begin{cousinweb}
The front piece is \textbf{mezzo/mezza}, \textquotedblleft{}half, middle,\textquotedblright{} from Latin \textbf{medius} (cousin of English \emph{medium, mid}, and of French \emph{mi-} in \emph{midi}):

\noindent
\begin{tabularx}{\linewidth}{@{}>{\raggedright\arraybackslash}X>{\raggedright\arraybackslash}X>{\raggedright\arraybackslash}X>{\raggedright\arraybackslash}X@{}}
\toprule
\textbf{Italian} & \textbf{=} & \textbf{$\leftarrow$ Latin} & \textbf{literally} \\
\midrule
\textbf{mezzogiorno} & \emph{mezzo} + \emph{giorno} & \emph{medius} + \emph{diurnum} & \textquotedblleft{}\textbf{mid}-\textbf{day}\textquotedblright{} (noon) \\
\textbf{mezzanotte} & \emph{mezza} + \emph{notte} & \emph{media} + \emph{noctem} & \textquotedblleft{}\textbf{mid}-\textbf{night}\textquotedblright{} (midnight) \\
\bottomrule
\end{tabularx}

\begin{itemize}
\raggedright
  \item \textbf{mezzogiorno} = \textquotedblleft{}mid-day.\textquotedblright{} \emph{giorno} (\textquotedblleft{}day\textquotedblright{}) is from Latin \emph{diurnum} (\textquotedblleft{}of the day\textquotedblright{}) — the same root as English \emph{journal} and \emph{journey} (a day's travel). So noon is \textquotedblleft{}the middle of the day-span.\textquotedblright{}
  \item \textbf{mezzanotte} = \textquotedblleft{}mid-night.\textquotedblright{} \emph{notte} (\textquotedblleft{}night\textquotedblright{}) $\leftarrow$ \emph{noctem}, twin of English \emph{night} and French \emph{nuit}. \emph{mezza} is the \textbf{feminine} \textquotedblleft{}half\textquotedblright{} (agreeing with the feminine \emph{notte}), where \emph{mezzo} is masculine (with \emph{giorno}) — the article system you just met, again.
\end{itemize}

To use them, they \emph{are} the hour:

\begin{quote}
\textbf{È mezzogiorno.} — \textquotedblleft{}It is noon.\textquotedblright{} \textbf{È mezzanotte.} — \textquotedblleft{}It is midnight.\textquotedblright{}
\end{quote}
\end{cousinweb}

\subsection*{Your turn}
Say these aloud:

\begin{itemize}
\raggedright
  \item \textquotedblleft{}mezzogiorno\textquotedblright{} (noon), \textquotedblleft{}mezzanotte\textquotedblright{} (midnight)
  \item break them — \textquotedblleft{}mezzo (middle) + giorno (day)\textquotedblright{}, \textquotedblleft{}mezza + notte (night)\textquotedblright{}
  \item \textquotedblleft{}È mezzogiorno. È mezzanotte.\textquotedblright{} — singular \emph{è}, no number
\end{itemize}

\begin{tcolorbox}[breakable,colback=teal!4,colframe=teal!35!black,title={Before you move on}]
What do \emph{mezzogiorno} and \emph{mezzanotte} literally mean? (\textquotedblleft{}Mid-day\textquotedblright{} and \textquotedblleft{}mid-night.\textquotedblright{}) What is \emph{mezzo/mezza} from, and its French cousin? (Latin \emph{medius}; French \emph{mi-} in \emph{midi/minuit}.) Why \emph{mezza}notte but \emph{mezzo}giorno? (\emph{mezza} is feminine for \emph{notte}, \emph{mezzo} masculine for \emph{giorno}.) Next chapter builds on this toward months and family.
\end{tcolorbox}
